% !TEX root = ../main.tex
\begin{table}[H]
    \centering
    \caption{Bilans debaty względem rundy niezależnej (liczba zadań, dla każdej debaty $n = 250$)}
    \label{tab:bilans_debaty}
    \small
    \setlength{\tabcolsep}{2.5pt}
    \renewcommand{\arraystretch}{1.35}
    \sbox0{%
    \begin{tabular}{|>{\centering\arraybackslash}p{2.0cm}|>{\centering\arraybackslash}p{1.4cm}|>{\centering\arraybackslash}p{1.3cm}|>{\centering\arraybackslash}p{1.9cm}|>{\centering\arraybackslash}p{1.3cm}|>{\centering\arraybackslash}p{1.9cm}|>{\centering\arraybackslash}p{1.3cm}|>{\centering\arraybackslash}p{1.9cm}|}
        \hline
        \rowcolor{gray!20}
        \multicolumn{2}{|c|}{} & \multicolumn{2}{c|}{\shortstack{\textbf{Zgoda}\\\textbf{od razu}}} & \multicolumn{2}{c|}{\shortstack{\textbf{Żaden model}\\\textbf{poprawny}}} & \multicolumn{2}{c|}{\shortstack{\textbf{Jeden model}\\\textbf{poprawny}}} \\ \cline{3-8}
        \rowcolor{gray!20}
        \textbf{Zbiór} & \textbf{Debata} & \textbf{razem} & \textbf{w tym poprawna} & \textbf{razem} & \textbf{w tym uratowane} & \textbf{razem} & \textbf{w tym utracone} \\ \hline
        Matematyka & Q+Q & 185 & 134 & 37 & 19 & 28 & 1 \\ \cline{2-8}
         & Q+G & 141 & 132 & 36 & 10 & 73 & 1 \\ \cline{2-8}
         & G+G & 207 & 185 & 30 & 8 & 13 & 3 \\ \hline
        Medycyna & Q+Q & 220 & 163 & 7 & 0 & 23 & 11 \\ \cline{2-8}
         & Q+G & 160 & 129 & 16 & 1 & 74 & 34 \\ \cline{2-8}
         & G+G & 229 & 150 & 10 & 1 & 11 & 6 \\ \hline
    \end{tabular}}
    \usebox{0}
    \par\smallskip
    \begin{minipage}{\wd0}
        \raggedright\footnotesize
        Zgoda od razu -- zadania, w których oba modele podały równoważne odpowiedzi w rundzie niezależnej, przed debatą; „w tym poprawna” -- gdy ta odpowiedź jest poprawna. Pozostałe zadania trafiły do debaty. Żaden model poprawny i Jeden model poprawny -- zadania z debatą, w których w rundzie niezależnej poprawną odpowiedź podał odpowiednio żaden albo jeden z modeli. Uratowane -- zadania z pierwszej z tych grup, w których wynik końcowy jest poprawny. Utracone -- zadania z drugiej, w których wynik końcowy nie jest poprawny, także gdy debata zakończyła się abstain.
    \end{minipage}
\end{table}
